\section{A comparator from a successor-seeded scalar recursor}
\label{sec:sr-comparator}

This section works in ordinary intuitionistic first-order logic with
actual equality. The numerical language has $0$ and the scalar recursor
\[
 R(n,a;v.h),\qquad R(0,a;v.h)=Sa,\qquad
 R(Sn,a;v.h)=h[R(n,a;v.h)/v],
\]
where $S a$ abbreviates $R(0,a;v.0)$. Only $v$ is bound in $h$;
every other parameter remains fixed during a recursion. Besides the
computation equations we assume zero/successor coverage,
$(0=1)\to\bot$ with $1=S0$, and contextual induction on a single literal
equation. Every induction below has such an equation as its invariant.
Case analysis on a previously proved finite range is an intuitionistic
logical elimination, not an induction on that range formula. We do not
assume predecessor, successor injectivity, ordinary initialized scalar
iteration, addition, pairing, or an equality decision function.

For each fixed update $f(v)$, the notation $R(n,a;f)$ means its fixed
finite term instance. Function notation below is abbreviation, never
quantification over numerical functions. $S^k a$ uses a fixed standard
number $k$ of successor signs.

\subsection{A positive selector and two finite cycles}

Define
\[
 T(n,a,b)=R(n,a;v.b),\qquad Z(n)=T(n,0,0).
\]
Their computation equations are
\[
 T(0,a,b)=Sa,\quad T(Sn,a,b)=b,\qquad Z(0)=1,\quad Z(Sn)=0.
\]
Coverage therefore proves that $Z$ takes only the values $0,1$.
Put
\[
 N(n)=Z(R(n,0;Z)),\qquad
 j(u)=T\bigl(u,0,T(Z(N(u)),1,0)\bigr),
 \qquad M(n)=j(j(R(n,0;j))).
\]
The range of $R(n,0;Z)$ is contained in $\{0,1\}$: its zero value is
$1$, and each successor value is a value of $Z$. Hence $N$ is Boolean.
Direct computation gives
\[
 N(0)=0,\qquad N(Sn)=Z(N(n)),\qquad
 N(S^2n)=N(n),\qquad N(S^3n)=Z(N(n)).                 \tag{N}
\]
The last two identities use only the Boolean range of $N$.

Coverage on $u$ and the Boolean range of $N(u)$ show that $j(u)$ is
always one of $0,1,2$. Finite computations give
\[
 j(0)=1,\qquad j(1)=2,\qquad j(2)=0.
\]
Consequently the range of $M$ is contained in $\{0,1,2\}$ and
\[
 M(0)=0,\quad M(Sn)=j(M(n)),\quad M(S^3n)=M(n),
 \qquad M(1)=1,\quad M(2)=2.                         \tag{M}
\]
Here $M(Sn)=j(M(n))$ follows by expanding one recursor step; the
three-period identity follows by three such steps and the finite cycle.
No induction asserting a disjunctive range has been used.

\subsection{The transient-seed construction}

For a fixed map $f$ and a fixed seed $a$, the composite
\[
                       J_f(n)=f(R(n,a;f))
\]
satisfies $J_f(0)=f(Sa)$ and $J_f(Sn)=f(J_f(n))$. Thus a state $Sa$
outside an intended orbit can be used once to enter that orbit. The
following constructions exploit a residue absent from the subsequent
orbit; they do not presuppose a general initialized iterator.

Define
\begin{align*}
 f_2(u)&=T(N(u),Su,0),&
 d_2(n)&=f_2(R(n,0;f_2)),\\
 f_3(u)&=T(M(u),S^2u,0),&
 d_3(n)&=f_3(R(n,0;f_3)).
\end{align*}
Since $f_2(1)=f_3(1)=0$, both $d_2(0)$ and $d_3(0)$ are zero.
Literal-equation induction proves respectively
\[
                     N(d_2(n))=0,\qquad M(d_3(n))=0.
\]
For the first step, the induction equation turns $f_2(d_2(n))$ into
$S^2d_2(n)$, whose $N$ value is still zero by (N). The second proof
uses $f_3(d_3(n))=S^3d_3(n)$ and (M) in the same way. We obtain the
entire recursion equations
\[
 d_2(0)=0,\quad d_2(Sn)=S^2d_2(n),\qquad
 d_3(0)=0,\quad d_3(Sn)=S^3d_3(n).                  \tag{Scale}
\]
These justify the intended doubling and tripling roles internally,
without using their standard-model meanings as premises. A separate
equation induction, using (N) and (Scale), proves
\[
                              N(d_3(n))=N(n).       \tag{Parity-scale}
\]
Indeed the successor step applies $Z$ to both sides of the current
induction equation.

Now put
\[
 f(u)=T\bigl(M(u),u,T(Z(N(M(u))),Su,0)\bigr),\qquad
 G(n)=f(R(n,1;f)).
\]
On the three possible residues of $M(u)$, finite computations yield
\[
 \begin{array}{c|ccc}
 M(u)&0&1&2\\ \hline
 f(u)&Su&S^2u&0.
 \end{array}                                                        \tag{Ramp-map}
\]
In particular $f(2)=0$. The seed $S1=2$ is therefore transient, and
\[
                         G(0)=0,\qquad G(Sn)=f(G(n)).
\]
Equation induction proves the central invariant
\[
                              M(G(n))=N(n).          \tag{Ramp-invariant}
\]
At the successor step retain this one induction equation and split the
Boolean $N(n)$. If it is zero, (Ramp-map) gives $G(Sn)=SG(n)$, and
the new $M$ value is $j(0)=1=N(Sn)$. If it is one, the new state is
$S^2G(n)$, and its $M$ value is $j^2(1)=0=N(Sn)$.
Thus
\[
 G(Sn)=
 \begin{cases}SG(n)&N(n)=0,\\S^2G(n)&N(n)=1,\end{cases}\qquad
 G(S^2n)=S^3G(n).                                    \tag{Ramp-step}
\]
One more step gives $G(S^3n)=S^4G(n)$ in the first Boolean case and
$G(S^3n)=S^5G(n)$ in the second.

Equation induction on $n$, using (Scale) and (Ramp-step), gives
\[
                              G(d_2(n))=d_3(n).       \tag{Ramp-double}
\]
For the other scaling identity define the conditional increment
\[
                    \Delta(e,a)=T(Z(e),a,a).
\]
It satisfies $\Delta(0,a)=a$ and $\Delta(1,a)=Sa$. We claim
\[
          G(d_3(n))=\Delta\bigl(N(n),d_3(G(n))\bigr). \tag{Ramp-triple}
\]
This is again one literal-equation induction on $n$. The base is zero.
For the step write $b=d_3(G(n))$ and use (Parity-scale). If $N(n)=0$,
the induction equation says $G(d_3(n))=b$. The new left side is
$G(S^3d_3(n))=S^4b$. The new right side is
$\Delta(1,d_3(SG(n)))=S(S^3b)=S^4b$.
If $N(n)=1$, the induction equation says $G(d_3(n))=Sb$.
The new left side is $S^5(Sb)=S^6b$, and the new right side is
$\Delta(0,d_3(S^2G(n)))=S^6b$. This proves the identity at every
argument, including nonstandard arguments in a source model.

\subsection{A factor map with small affine constants}

Define
\[
 F_2(u)=S d_2(u),\qquad F_3(u)=S^2d_3(u),\qquad
 H(u)=T\bigl(Z(N(u)),G(u),u\bigr).
\]
The preceding equations give
\[
 N(F_2(u))=1,\qquad N(F_3(u))=N(u),\qquad
 H(u)=\begin{cases}u&N(u)=0,\\SG(u)&N(u)=1.\end{cases}             \tag{Factor-cases}
\]
We now prove
\[
                H(F_2(u))=F_3(u),\qquad
                H(F_3(u))=F_3(H(u)).                  \tag{Factor}
\]
For the first, $N(d_2(u))=0$ and (Ramp-double) give
\[
 H(F_2(u))=S G(Sd_2(u))=S^2G(d_2(u))=S^2d_3(u).
\]
For the second, split the Boolean $N(u)$. At zero both sides are
$F_3(u)$. At one, (Ramp-step) and (Ramp-triple) give
\[
 H(F_3(u))=S G(S^2d_3(u))=S^4G(d_3(u))=S^5d_3(G(u)),
\]
whereas $F_3(H(u))=S^2d_3(SG(u))=S^5d_3(G(u))$ by (Scale).
Notice also that $H(2)=2$, since $N(2)=0$.

\subsection{Removing affine factors gives comparison}

Put
\[
 E(x)=R(x,0;v.S^2d_2(v)),\qquad Q(a)=S^4d_3(a).
\]
Their useful equations are
\[
 E(0)=1,\qquad S E(Sx)=F_2(S E(x)),\qquad S Q(a)=F_3(Sa).          \tag{Seeds}
\]
The middle equation follows from $d_2(Sa)=S^2d_2(a)$; the last from
$d_3(Sa)=S^3d_3(a)$.

For a fixed update $h$, if $h(Sa)=Sb$, equation induction on $y$
proves
\[
                         R(Sy,a;h)=R(y,b;h).          \tag{Rotation}
\]
Its base is the assumed equation, and its step applies $h$ to the
current induction equation. All initializer and update parameters stay
fixed. In particular, (Factor) and (Seeds) supply
\[
 H(S E(Sx))=F_3(S E(x))=S Q(E(x)),
\]
so (Rotation) gives
\[
                 R(Sy,E(Sx);H)=R(y,Q(E(x));H).        \tag{Seed-rotation}
\]
A further equation induction on $y$, with $a$ fixed, proves
\[
                  R(y,Q(a);H)=F_3(R(y,a;H)).          \tag{Iterator-factor}
\]
Its base is the last equation of (Seeds). Its step is precisely the
second equation of (Factor), applied after the induction equation.

Finally define the finite scalar term
\[
                    L(x,y)=Z\bigl(N(R(y,E(x);H))\bigr).
\]
Its Boolean range is immediate from that of $N$ and $Z$. The complete
comparison recursion specification is internally provable:
\[
 \boxed{L(0,y)=1,\qquad L(Sx,0)=0,\qquad L(Sx,Sy)=L(x,y).}         \tag{Compare}
\]
For the first equation, induction on the literal equation
$R(y,1;H)=2$ uses $H(2)=2$, and then $Z(N(2))=1$.
For the second, (Seeds) makes the initial state $F_2(S E(x))$, whose
$N$ value is one. For the third, (Seed-rotation), (Iterator-factor), and
$N(F_3(u))=N(u)$ give the result directly.

We will also need persistence:
\[
                          L(x,y)=1\ \longrightarrow\ L(x,Sy)=1. \tag{Persist}
\]
Let $r=R(y,E(x);H)$. The Boolean range of $N(r)$ and $0\ne1$ show
from $Z(N(r))=1$ that $N(r)=0$. Thus $H(r)=r$ by (Factor-cases),
so the next recursor state is unchanged. No induction on an implication
occurs here. Two ordinary equation inductions applied to (Compare)
also give $L(n,n)=1$ and $L(Sn,n)=0$.

Every construction above is a fixed finite term in $0,R$. The intended
standard meaning $L(x,y)=1$ iff $x\le y$ explains the design, but neither
that order interpretation nor evaluation only at standard numerals is
used in the proof. This completes the comparator bootstrap before
predecessor, an unrestricted selector, or successor injectivity.
